\section{Analyse} 
\label{sec:Analyse}
Im Folgenden wird zunächst die bisherige Situation im Unternehmen betrachtet.
Darauf aufbauend wird untersucht, welchen wirtschaftlichen Nutzen das Projekt bietet.
Anschließend werden die fachlichen und technischen Anforderungen an die Anwendung beschrieben.
Zum Abschluss folgt eine Nutzwertanalyse, um den nicht-monetären Nutzen des Projektes zu bewerten.

\subsection{Ist-Analyse} 
\label{sec:IstAnalyse}
Aktuell werden kundenbezogene Daten in zwei voneinander getrennten Systemen gespeichert.
Zum einen im \textit{user-service}, welcher Daten für das Kundenportal verwaltet, und zum anderen im CRM-System \textit{HubSpot}.
Zwischen beiden Systemen besteht bislang keine eigenständig verwaltete und durchgängig nachvollziehbare Synchronisation.
Dadurch kann es passieren, dass Kontaktinformationen eines Kunden in beiden Systemen unterschiedlich vorliegen.
\par\smallskip
Besonders problematisch ist dies bei Informationen, die für die Kommunikation mit dem Kunden relevant sind.
Dazu zählen unter anderem die E-Mail-Adresse, der Vorname, der Nachname, die Anrede sowie Kommunikationspräferenzen.
Werden diese Daten in einem der beiden Systeme geändert, ist nicht sichergestellt, dass die Änderung zuverlässig im jeweils anderen System übernommen wird.
\par\smallskip
Zusätzlich existiert bislang keine zentrale Stelle, die nachvollziehbar speichert, wann eine Synchronisation durchgeführt wurde,
aus welchem System eine Änderung stammt und ob diese erfolgreich verarbeitet wurde.
Fehler bei der Übertragung könnten dadurch nur mit erhöhtem manuellem Aufwand erkannt und korrigiert werden.
Ebenso gibt es bisher keine eigenständige technische Lösung, um mit temporären Fehlern, Systemausfällen oder Rate Limits der HubSpot API umzugehen.
\par\smallskip
Aus diesem Grund soll mit dem Projekt \textit{Nightcrawler} ein eigenständiger Microservice entwickelt werden,
der die Synchronisation zwischen beiden Systemen übernimmt.
Der Name ist an den teleportierenden Helden angelehnt, da die Anwendung Kundendaten kontrolliert zwischen zwei voneinander getrennten Systemen überträgt.
Hierzu soll eine PostgreSQL-Datenbank eingeführt werden, in der sowohl der aktuelle Synchronisationszustand
zwischen Benutzer und Kontakt als auch die Historie aller Synchronisationsereignisse gespeichert wird \cite{PostgreSQL16Docs}.
\par\smallskip
Die Kommunikation zwischen den beteiligten Systemen soll asynchron über RabbitMQ erfolgen.
Die bisherige Situation wird dadurch verbessert, dass Kundendaten automatisch, nachvollziehbar und fehlertolerant synchronisiert werden können.
Zusätzlich soll die neue Anwendung Konflikte bei gleichzeitigen Änderungen erkennen und nach festgelegten Regeln auflösen.


\subsection{Wirtschaftlichkeitsanalyse}
\label{sec:Wirtschaftlichkeitsanalyse}
Für die Umsetzung des Projektes sind insgesamt 80 Stunden vorgesehen.
Bei einem internen Verrechnungssatz von 35 Euro Netto pro Stunde ergeben sich folgende Entwicklungskosten:

\[
80 \text{ Stunden} \times 35 \text{ Euro/Stunde} = 2800 \text{ Euro}
\]

Dem gegenüber steht der Nutzen der automatisierten Synchronisation.
Ohne eine solche Lösung müssten Abweichungen zwischen den Datenbeständen manuell erkannt, überprüft und korrigiert werden.
Dies würde insbesondere im Kundensupport und in der technischen Betreuung zusätzlichen Aufwand verursachen.
\par\smallskip
Wird angenommen, dass durch die automatisierte Synchronisation wöchentlich 2,5 Stunden manueller Aufwand eingespart werden können,
ergibt sich bei einem internen Stundensatz von 15,50 Euro Netto pro Stunde ein jährliches Einsparpotenzial von:

\[
2{,}5 \text{ Stunden/Woche} \times 15{,}50 \text{ Euro/Stunde} \times 52 \text{ Wochen} = 2015 \text{ Euro/Jahr}
\]

Die Amortisationsdauer der Eigenentwicklung beträgt damit:

\[
2800 \text{ Euro} \div 2015 \text{ Euro/Jahr} \approx 1{,}39 \text{ Jahre}
\]

Die Entwicklungskosten sind somit nach ungefähr 17 Monaten durch die eingesparte Arbeitszeit ausgeglichen.
\par\smallskip
Im Rahmen der Make-or-Buy-Entscheidung wurde außerdem geprüft, ob eine bestehende Integrationsplattform oder ein fertiger HubSpot-Connector eingesetzt werden kann.
Eine solche Buy-Lösung hätte zwar einen geringeren Entwicklungsaufwand, würde jedoch laufende Lizenzkosten verursachen und könnte die projektspezifischen Anforderungen nur eingeschränkt erfüllen.
Dazu zählen insbesondere die eigene Synchronisationshistorie, die definierte Konfliktbehandlung, die Retry-Logik sowie die enge Integration in die bestehende Systemarchitektur.
\par\smallskip
Für eine Buy-Lösung wird beispielhaft ein einmaliger Einrichtungsaufwand von 16 Stunden sowie ein monatlicher Lizenzbetrag von 250 Euro Netto angesetzt:

\[
16 \text{ Stunden} \times 35 \text{ Euro/Stunde} + 250 \text{ Euro/Monat} \times 12 \text{ Monate} = 3560 \text{ Euro}
\]

Damit liegen die Kosten der Buy-Lösung bereits im ersten Jahr über den Entwicklungskosten der Eigenentwicklung.
Zusätzlich bietet die Eigenentwicklung einen höheren fachlichen Nutzen und eine bessere Erweiterbarkeit.
Aus wirtschaftlicher und fachlicher Sicht wird daher die Eigenentwicklung des \textit{Nightcrawlers} umgesetzt.


\subsection{Anforderungsanalyse}
\label{sec:Anforderungsanalyse}
Damit das Projektziel erreicht werden kann, muss die Anwendung sowohl fachliche als auch technische Anforderungen erfüllen.
\par\smallskip
Fachlich muss der \textit{Nightcrawler} in der Lage sein, Änderungen an Kundendaten zwischen dem \textit{user-service} und \textit{HubSpot} in beide Richtungen zu synchronisieren.
Dabei sollen mindestens die Datenfelder \textit{email}, \textit{givenName}, \textit{surname}, \textit{salutation} und \textit{CommunicationPreference} berücksichtigt werden.
\par\smallskip
Wenn ein Benutzer im \textit{user-service} erstellt oder geändert wird, soll ein entsprechendes Ereignis verarbeitet und die dazugehörigen Daten in \textit{HubSpot} aktualisiert werden.
Wird ein Benutzer neu erstellt, soll zunächst geprüft werden, ob anhand der E-Mail-Adresse bereits ein passender Kontakt in \textit{HubSpot} existiert.
Ist dies der Fall, soll dieser Kontakt aktualisiert werden.
Falls kein Kontakt gefunden wird, soll ein neuer Kontakt erstellt werden.
\par\smallskip
Wenn dagegen in \textit{HubSpot} ein Kontakt geändert wird, soll diese Änderung ebenfalls verarbeitet und an den \textit{user-service} weitergegeben werden.
Kontakte aus \textit{HubSpot}, für die kein passender Benutzer im \textit{user-service} existiert, sollen zunächst ignoriert werden.
\par\smallskip
Damit jede Synchronisation nachvollziehbar bleibt, muss jedes verarbeitete Ereignis zunächst in einer Historie gespeichert werden.
Diese \textit{SyncHistory} soll Informationen zur Quelle des Ereignisses, zum Ereignistyp, zu Benutzer- und Kontaktzuordnung,
zum Payload, zum Zeitpunkt des Auftretens, zum aktuellen Status sowie zur Anzahl der bisherigen Verarbeitungsversuche enthalten.
\par\smallskip
Zusätzlich muss ein \textit{UserContactState} gespeichert werden.
Dieser dient dazu, die Zuordnung zwischen Benutzer und HubSpot-Kontakt sowie den letzten bekannten Synchronisationsstand zu verwalten.
Dabei sollen unter anderem die Zeitpunkte der letzten Änderung im \textit{user-service}, der letzten Änderung in \textit{HubSpot}
und der letzten erfolgreichen Synchronisation gespeichert werden.
\par\smallskip
Technisch muss die Anwendung acht Event Queues über RabbitMQ verwenden.
Da keine Nachrichten bei einem Neustart oder Ausfall verloren gehen sollen, müssen die Hauptqueues als durable Quorum Queues umgesetzt werden.
Für temporäre Fehler soll eine Retry-Logik mit TTL-basierten Retry Queues eingesetzt werden.
Nicht behandelbare Fehler sollen in einer Dead-Letter-Queue abgelegt werden \cite{RabbitMQDocs}.
\par\smallskip
Zur Konfliktbehandlung soll eine optimistische Strategie verwendet werden.
Wenn Daten in beiden Systemen unterschiedlich geändert wurden, soll der Datenstand des \textit{user-service} bevorzugt werden.
Wird also eine Änderung aus \textit{HubSpot} erkannt, die im Konflikt mit aktuelleren Daten aus dem \textit{user-service} steht,
soll diese Änderung nicht übernommen, sondern stattdessen eine neue Synchronisation in Richtung \textit{HubSpot} angestoßen werden.

\subsection{Nutzwertanalyse}
\label{sec:Nutzwertanalyse}
Neben dem wirtschaftlichen Nutzen bietet das Projekt auch einen nicht-monetären Mehrwert.
Zur Bewertung werden die Kriterien Datenkonsistenz, Nachvollziehbarkeit, Ausfallsicherheit, Wartbarkeit und Automatisierungsgrad betrachtet.
Dabei werden die bisherige Ist-Situation, eine mögliche Buy-Lösung und die Eigenentwicklung des \textit{Nightcrawlers} verglichen.

\tabelle{Nutzwertvergleich}{tab:Nutzwert}{Nutzwert.tex}

Aus der Nutzwertanalyse ergibt sich, dass eine Buy-Lösung gegenüber der bisherigen Situation zwar eine Verbesserung darstellt,
jedoch nicht alle projektspezifischen Anforderungen vollständig erfüllt.
Die Eigenentwicklung erreicht den höchsten Nutzwert, da sie besonders bei Nachvollziehbarkeit, Konfliktbehandlung,
Ausfallsicherheit und Anpassbarkeit an die bestehende Systemarchitektur Vorteile bietet.
Die Entscheidung zugunsten des \textit{Nightcrawlers} wird dadurch zusätzlich bestätigt.